% 第 1 章 Introduction to PrimeTime
\bichapter{PrimeTime 简介}{Introduction to PrimeTime}
\label{chap:intro}

PrimeTime Suite 执行全芯片、门级静态时序分析（static timing analysis，STA），这是芯片设计与分析流程中的关键环节。该工具通过检查所有路径是否存在时序违例，对设计的时序性能进行穷尽验证，无需逻辑仿真或测试向量。要了解 PrimeTime 工具的基础内容，请参阅：
\begin{itemize}
  \item PrimeTime 能力
  \item PrimeTime 产品层级与许可证
  \item PrimeTime 附加工具
  \item 在实现流程中使用 PrimeTime
  \item 与 Synopsys 实现工具的兼容性
  \item 静态时序分析概览
\end{itemize}

% ============================================================
\section{PrimeTime 能力}
\subsection*{PrimeTime Capabilities}

在 PrimeTime 中，可执行多种设计检查与分析：
\begin{itemize}
  \item \textbf{设计检查（Design Checks）}
  \begin{itemize}
    \item Setup、hold、recovery 与 removal 约束
    \item 用户指定的 data-to-data 时序约束
    \item 时钟门控（clock-gating）的 setup 与 hold 约束
    \item 时钟的最小周期（minimum period）与最小脉冲宽度（minimum pulse width）
    \item 设计规则（最小/最大 transition time、capacitance 与 fanout）
  \end{itemize}
  \item \textbf{分析功能（Analysis Features）}
  \begin{itemize}
    \item 多时钟与多时钟频率
    \item 伪路径（false path）与多周期路径（multicycle path）时序例外
    \item 透明锁存器（transparent latch）分析与时间借用（time borrowing）
    \item 对 setup 与 hold 约束同时进行最小与最大延时分析
    \item 考虑工艺、电压与温度（PVT）片上变化（on-chip variation，OCV）的分析
    \item 对指定输入施加常量或特定跳变的 case analysis
    \item 按模块特定工作模式（如 RAM 模块的读模式或写模式）进行的 mode analysis
    \item 瓶颈分析（bottleneck analysis），报告导致最多时序违例的单元
    \item 物理相邻网络之间的串扰（crosstalk）分析
    \item 使用 HyperScale 技术的高效层次化分析
    \item 全面的约束一致性检查（constraint consistency checking）
    \item 工程变更指令（ECO）分析：通过插入缓冲器、调整单元尺寸与交换单元来纠正违例并优化功耗
  \end{itemize}
\end{itemize}

% ============================================================
\section{PrimeTime 产品层级与许可证}
\subsection*{PrimeTime Product Tiers and Licenses}

PrimeTime 产品提供三个层级：PrimeTime Base、PrimeTime Elite 与 PrimeTime Apex。

表~\ref{tab:pt-tiers-licenses} 列出各产品层级所包含的许可证。较新功能使用 PrimeTime-ELT 与 PrimeTime-APX 许可证；较旧功能使用其他许可证。

\begin{table}[htbp]
\centering
\caption{PrimeTime 产品打包层级所含许可证}
\label{tab:pt-tiers-licenses}
\begin{tabular}{|l|l|l|}
\hline
\textbf{PrimeTime Base} & \textbf{PrimeTime Elite} & \textbf{PrimeTime Apex} \\
\hline
PrimeTime & PrimeTime & PrimeTime $\times$ 2 \\
PrimeTime-SI & PrimeTime-SI & PrimeTime-SI $\times$ 2 \\
 & PrimeTime-ADV & PrimeTime-ADV $\times$ 2 \\
 & PrimeTime-ADV-PLUS & PrimeTime-ADV-PLUS $\times$ 2 \\
 & PrimeTime-ELT & PrimeTime-ELT $\times$ 2 \\
 &  & PrimeTime-APX \\
\hline
\end{tabular}
\end{table}

表~\ref{tab:pt-tier-features} 列出各 PrimeTime 产品层级所含功能，以及启用该功能所需的许可证。

\begin{table}[htbp]
\centering
\caption{PrimeTime 产品层级与所含功能}
\label{tab:pt-tier-features}
\small
\begin{tabularx}{\textwidth}{|X|c|c|c|}
\hline
\tblhead{功能（及许可证）} & \tblhead{Base} & \tblhead{Elite} & \tblhead{Apex} \\
\hline
多核处理（Multicore processing） & X（16 核） & X（32 核） & X（64 核） \\
信号完整性分析（PrimeTime-SI） & X & X & X \\
静态噪声分析（PrimeTime-SI） & X & X & X \\
约束一致性（PrimeTime-SI） & X & X & X \\
以二进制格式写入与读取 Context（PrimeTime-SI） & X & X & X \\
IR Drop Annotation（PrimeTime-SI） & X & X & X \\
用 \cmd{write\_spice\_deck} 生成 SPICE Deck（PrimeTime-SI） & X & X & X \\
参数化片上变化 POCV（PrimeTime-ADV） &  & X & X \\
多输入开关分析 Multi-Input Switching（PrimeTime-ADV） &  & X & X \\
物理感知 ECO（PrimeTime-ADV） &  & X & X \\
基于核的分布式许可证 Distributed Core-Based Licensing（PrimeTime-ADV） &  & X & X \\
在 PrimeTime 中回放 ECO Change List（PrimeTime-ADV） &  & X & X \\
基于机器学习的穷尽路径分析（PrimeTime-ADV-PLUS） &  & X & X \\
HyperTrace 加速路径分析（PrimeTime-ADV-PLUS） &  & X & X \\
同时多电压分析 SMVA（PrimeTime-ADV-PLUS） &  & X & X \\
基于矩建模的 POCV（PrimeTime-ADV-PLUS） &  & X & X \\
含 Via Variation 的 POCV（PrimeTime-ADV-PLUS） &  & X & X \\
多输入开关分析——\cmd{lib\_arc} 与高级模式（PrimeTime-ADV-PLUS） &  & X & X \\
基于机器学习的加速功耗回收（PrimeTime-ADV-PLUS） &  & X & X \\
GUI 中的 Layout View 与 ECO（PrimeTime-ADV-PLUS） &  & X & X \\
GPD Parasitic Explorer（PrimeTime-ADV-PLUS） &  & X & X \\
将 Ansys RedHawk-SC 用于时序分析（PrimeTime-ADV-PLUS） &  & X & X \\
HyperGrid 分布式分析（PrimeTime-ADV-PLUS） &  & X & X \\
高级数据管理 Advanced Data Management（PrimeTime-ELT） &  & X & X \\
Context Budgeting（PrimeTime-APX） &  &  & X \\
\hline
\end{tabularx}
\end{table}

% ============================================================
\section{PrimeTime 附加工具}
\subsection*{PrimeTime Add-On Tools}

PrimeTime Suite 支持以下附加功能：
\begin{itemize}
  \item \textbf{PrimeECO}——支持工程变更指令（ECO）实现与分析：
  \begin{itemize}
    \item PrimeECO 附加工具将 PrimeTime、IC Compiler II 与 StarRC 的时序分析、物理实现与寄生提取功能整合为单一集成工具。它可执行多轮 ECO 修复迭代，无需在迭代之间暂停。更多信息见 \textit{PrimeECO User Guide}。
  \end{itemize}
  \item \textbf{PrimePower}——支持高级功耗分析：
  \begin{itemize}
    \item PrimePower 附加工具分析基于单元设计的全芯片功耗。支持无向量与基于向量的峰值功耗及平均功耗分析。更多信息见 \textit{PrimePower User Guide}。
  \end{itemize}
  \item \textbf{PrimeShield}——支持参数化鲁棒性分析：
  \begin{itemize}
    \item 在先进工艺节点下，面对日益增大的工艺与电压变异性，PrimeShield 提供设计鲁棒性分析与修复，帮助设计者减少悲观、过度裕量与过度设计，从而有效降低功耗并提升最高工作频率。更多信息见 \textit{PrimeShield User Guide}。
  \end{itemize}
\end{itemize}

% ============================================================
\section{在实现流程中使用 PrimeTime}
\subsection*{Using PrimeTime in the Implementation Flow}

PrimeTime 工具可与 Design Compiler、IC Compiler II、Fusion Compiler 与 StarRC 等其他 Synopsys 工具良好协同。这些工具共享许多相同的库、数据库与命令，如下图所示。

\figplaceholder{Figure 1: Implementation Flow With PrimeTime Analysis}{含 PrimeTime 分析的实现流程}{fig:impl-flow}

从 RTL 设计出发，Design Compiler 生成门级设计。基于门级网表与物理库信息，物理实现工具（IC Compiler II）执行布局与布线。StarRC 寄生提取工具从芯片版图数据库提取寄生 RC 数据。

PrimeTime 读取门级网表与寄生数据，并利用逻辑库（\texttt{.db}）中的信息验证设计时序。若发现违例，可生成工程变更指令（ECO），指导物理实现工具修复违例并优化功耗。

PrimeTime 也可在其他设计流程中作为独立静态时序分析器运行。

% ============================================================
\section{与 Synopsys 实现工具的兼容性}
\subsection*{Compatibility With Synopsys Implementation Tools}

PrimeTime 静态时序分析工具旨在与 Design Compiler 综合工具以及 IC Compiler II 布局布线工具良好协同。兼容性体现在：
\begin{itemize}
  \item 使用相同的逻辑库，并读取相同的设计数据文件；
  \item 支持 Synopsys Design Constraints（SDC）格式以指定设计约束（含时序与面积约束）；
  \item 共享许多命令（如 \cmd{create\_clock}、\cmd{set\_input\_delay} 与 \cmd{report\_timing}），操作相同或相似；
  \item 共享相同的延时计算算法，产生相似的延时结果；
  \item 生成相似的时序报告；
  \item PrimeTime 可通过生成 change list 向 IC Compiler II 提供 ECO 指导，以修复时序违例并优化功耗。
\end{itemize}

尽管 Design Compiler 与 IC Compiler II 自身具备内置静态时序分析能力，PrimeTime 在速度、精度、容量与灵活性方面更优，并提供其他工具所不支持的许多功能。

\begin{seeAlsoBox}
在 Design Compiler 中使用 PrimeTime，见第~\ref{chap:pt-dc}~章。
\end{seeAlsoBox}

% ============================================================
\section{静态时序分析概览}
\subsection*{Overview of Static Timing Analysis}

静态时序分析是一种通过检查所有可能路径是否存在时序违例，来验证设计时序性能的方法。PrimeTime 将设计分解为时序路径，计算每条路径上的信号传播延时，并检查设计内部以及输入/输出接口处的时序约束违例。

另一种时序分析方法是动态仿真（dynamic simulation）：对给定输入激励向量集合，确定电路的完整行为。与动态仿真相比，静态时序分析快得多，因为不必仿真电路的逻辑操作；也更彻底，因为它检查所有时序路径，而不仅是特定测试向量所敏化的逻辑条件。但静态时序分析只能检查电路设计的时序，不能检查功能。

\subsection{时序路径}
\subsubsection*{Timing Paths}

进行时序分析时，PrimeTime 首先将设计分解为时序路径。每条时序路径由以下元素组成：
\begin{itemize}
  \item \textbf{Startpoint（起点）}——时序路径起点，数据由时钟边沿发射，或数据必须在特定时刻可用。每个起点必须是输入端口或寄存器时钟引脚。
  \item \textbf{Combinational logic network（组合逻辑网络）}——无存储或内部状态的元件。可包含 AND、OR、XOR 与反相器等，但不能包含触发器、锁存器、寄存器或 RAM。
  \item \textbf{Endpoint（终点）}——时序路径终点，数据由时钟边沿捕获，或数据必须在特定时刻可用。每个终点必须是寄存器数据输入引脚或输出端口。
\end{itemize}

下图给出简单设计示例中的时序路径。

\figplaceholder{Figure 2: Timing paths}{时序路径}{fig:timing-paths}

示例中，每个逻辑云表示一个组合逻辑网络。每条路径从数据发射点出发，经过若干组合逻辑，止于数据捕获点。

\begin{table}[htbp]
\centering
\caption{时序路径示例}
\begin{tabularx}{\textwidth}{|l|X|X|}
\hline
\tblhead{路径} & \tblhead{起点} & \tblhead{终点} \\
\hline
Path 1 & 输入端口 & 时序元件的数据输入 \\
Path 2 & 时序元件的时钟引脚 & 时序元件的数据输入 \\
Path 3 & 时序元件的时钟引脚 & 输出端口 \\
Path 4 & 输入端口 & 输出端口 \\
\hline
\end{tabularx}
\end{table}

一个组合逻辑云可能包含多条路径，如下图所示。PrimeTime 用最长路径计算最大延时，用最短路径计算最小延时。

\figplaceholder{Figure 3: Multiple paths through combinational logic}{组合逻辑中的多条路径}{fig:multi-paths}

PrimeTime 还考虑以下类型的路径用于时序分析：
\begin{itemize}
  \item \textbf{Clock path（时钟路径）}——从时钟输入端口或单元引脚出发，经过一个或多个缓冲器或反相器，到达时序元件时钟引脚的路径；用于数据 setup 与 hold 检查。
  \item \textbf{Clock-gating path（时钟门控路径）}——从输入端口到时钟门控元件的路径；用于时钟门控 setup 与 hold 检查。
  \item \textbf{Asynchronous path（异步路径）}——从输入端口到时序元件异步置位或清零引脚的路径；用于 recovery 与 removal 检查。
\end{itemize}

\figplaceholder{Figure 4: Types of paths considered for timing analysis}{时序分析所考虑的路径类型}{fig:path-types}

\subsection{延时计算}
\subsubsection*{Delay Calculation}

将设计分解为时序路径集合后，PrimeTime 计算每条路径上的延时。路径总延时为路径中所有单元延时与网络延时之和。

\subsubsection{单元延时}
\paragraph*{Cell Delay}

单元延时是路径中逻辑门从输入到输出的延时量。在没有来自 SDF 文件的反标延时信息时，工具根据该单元逻辑库中的延时表计算单元延时。

通常，延时表将延时量列为一个或多个变量（如输入 transition time 与输出负载电容）的函数。工具根据这些表项计算各单元延时。必要时，PrimeTime 对表值进行插值或外推，以得到当前设计条件下指定的延时值。

\subsubsection{网络延时}
\paragraph*{Net Delay}

网络延时是时序路径中从一个单元输出到下一单元输入的延时量。该延时由两单元之间互连的寄生电容、网络电阻以及驱动该网络的单元有限驱动能力共同引起。

PrimeTime 用以下方法计算网络延时：
\begin{itemize}
  \item 根据线负载模型（wire load model）估算延时——布局前芯片拓扑未知时使用；
  \item 使用从 Standard Delay Format（SDF）文件反标的具体时间值；
  \item 使用从 Galaxy Parasitic Database（GPD）、Standard Parasitic Exchange Format（SPEF）、Detailed Standard Parasitic Format（DSPF）或 Reduced Standard Parasitic Format（RSPF）文件反标的详细寄生电阻与电容数据。
\end{itemize}

\subsection{约束检查}
\subsubsection*{Constraint Checks}

确定时序路径并计算路径延时后，PrimeTime 检查时序约束违例，例如 setup 与 hold 约束：
\begin{itemize}
  \item \textbf{Setup 约束}规定：在捕获数据的时钟边沿之前，数据必须提前多少时间在时序器件输入端可用。该约束对数据路径相对于时钟边沿施加最大延时限制。
  \item \textbf{Hold 约束}规定：在捕获数据的时钟边沿之后，数据必须在时序器件输入端保持稳定多久。该约束对数据路径相对于时钟边沿施加最小延时限制。
\end{itemize}

除 setup 与 hold 外，PrimeTime 还可检查 recovery 与 removal 约束、data-to-data 约束、时钟门控 setup 与 hold 约束，以及时钟信号的最小脉冲宽度。

避免违例的时间裕量称为 \textbf{slack}（时序裕量）。例如，对 setup 约束，若信号须不晚于 8~ns 到达单元输入，而实际到达时间为 5~ns，则 slack 为 3~ns。Slack 为 0 表示约束刚好满足；负 slack 表示存在时序违例。

\subsubsection{触发器的 Setup 与 Hold 检查示例}
\paragraph*{Example of Setup and Hold Checks for Flip-Flops}

下例说明 PrimeTime 如何对触发器检查 setup 与 hold 约束。

\figplaceholder{Figure 5: Setup and hold checks}{Setup 与 hold 检查}{fig:setup-hold-ff}

本例假设逻辑库中定义触发器最小 setup 时间为 1.0 时间单位，最小 hold 时间为 0.0 时间单位。工具中定义时钟周期为 10 时间单位。时间单位大小（如 ns 或 ps）在逻辑库中指定。

默认情况下，工具假定信号在一个时钟周期内沿每条数据路径传播。因此，执行 setup 检查时，工具验证从 FF1 发射的数据在一个时钟周期内到达 FF2，并至少提前 1.0 时间单位到达，以便被 FF2 的下一时钟边沿捕获。若数据路径延时过长，则报告为时序违例。对该 setup 检查，工具考虑数据路径上可能的最长延时，以及 FF1 与 FF2 之间时钟路径上可能的最短延时。

执行 hold 检查时，工具验证从 FF1 发射的数据不早于上一时钟周期的捕获时钟边沿到达 FF2。这确保已存在于 FF2 输入端的数据在捕获上一周期数据的时钟边沿之后保持足够稳定。对该 hold 检查，工具考虑数据路径上可能的最短延时，以及 FF1 与 FF2 之间时钟路径上可能的最长延时。若时钟路径延时过长，可能发生 hold 违例。

\subsubsection{锁存器的 Setup 与 Hold 检查示例}
\paragraph*{Example of Setup and Hold Checks for Latches}

基于锁存器的设计通常使用两相、非重叠时钟控制数据路径中的相继寄存器。此时，PrimeTime 可利用时间借用（time borrowing）减轻对相继路径的约束。

例如，考虑下图所示的两相、基于锁存器的路径。三个锁存器均为电平敏感，当 G 输入为高时门控有效。L1 与 L3 由 PH1 控制，L2 由 PH2 控制。上升沿从锁存器输出发射数据，下降沿在锁存器输入捕获数据。本例中假设锁存器 setup 与延时时间为零。

\figplaceholder{Figure 6: Latch-based paths}{基于锁存器的路径}{fig:latch-paths}

对从 L1 到 L2 的路径，PH1 上升沿发射数据。数据必须在时间 = 20 的 PH2 下降沿之前到达 L2。该时序要求标记为 Setup~1。取决于 L1 与 L2 之间的延时量，数据可能在时间 = 10 的 PH2 上升沿之前或之后到达，时序图中用虚线箭头 ``Arrival a'' 与 ``Arrival b'' 表示。在时间 = 20 之后到达即为时序违例。

若数据在时间 = 10 的 PH2 上升沿之前到达 L2（Arrival a），则下一条从 L2 到 L3 的路径数据由时间 = 10 的 PH2 上升沿发射，如同同步触发器。此时不从第二条路径借用时间。L2 到 L3 的该时序要求标记为 Setup~2a。

若数据在 PH2 上升沿之后到达（Arrival b），则第一条路径（L1 到 L2）从第二条路径（L2 到 L3）借用时间。此时第二条路径的数据发射不发生在上升沿，而发生在 L2 的数据到达时刻，介于 PH2 上升沿与下降沿之间。该时序要求标记为 Setup~2b。发生借用时，路径起点为 L2 的 D 引脚而非 G 引脚。

对第一条路径（L1 到 L2），若发生借用，PrimeTime 报告 setup slack 为零；若数据在时间 = 10 的上升沿之前到达，slack 为正；若在时间 = 20 的下降沿之后到达，slack 为负（违例）。

执行 hold 检查时，PrimeTime 相对于 setup 检查考虑发射与捕获边沿，如下图所示。它验证起点发射的数据不会过快到达终点，从而确保上一周期发射的数据已被锁存，不会被新数据覆盖。

\figplaceholder{Figure 7: Hold checks in latch-based paths}{基于锁存器路径中的 hold 检查}{fig:latch-hold}

\subsection{时序例外}
\subsubsection*{Timing Exceptions}

若某些路径并非按 PrimeTime 默认假定的 setup/hold 行为工作，需将这些路径指定为时序例外（timing exceptions）；否则工具可能错误地将这些路径报告为时序违例。

PrimeTime 允许指定以下类型的例外：
\begin{itemize}
  \item \textbf{False path（伪路径）}——因逻辑配置、预期数据序列或工作模式而永不被敏化的路径。
  \item \textbf{Multicycle path（多周期路径）}——设计为从发射到捕获需要超过一个时钟周期的路径。
  \item \textbf{Minimum or maximum delay path（最小/最大延时路径）}——必须满足用户以时间值显式指定的延时约束的路径。
\end{itemize}

更多信息见第~\ref{chap:paths}~章「时序路径与例外」。
